% GENERATED FILE. Edit canonical lessons, then run npm run generate:books.
% canonical-source-hash: fnv1a64:5612ae1c66e920a6
% canonical-lessons: UR-C465-haddi, UR-C465-pasli, UR-C465-khopri, UR-C465-rirh, UR-C465-jabra

\chapter{Bones}
\label{ch:ur-bones}
\hlchaptermodality{465}

\begin{chapteropening}
\textbf{By the end of this chapter:} \emph{I can name some of the bones.}

Complete the last lesson of chapter 465: I can name some of the bones.
\end{chapteropening}

\section[\texorpdfstring{\ur{ہڈی} (haḍḍī)}{ہڈی (haḍḍī)}]{\ur{ہڈی} (haḍḍī) — a bone}

\label{lesson:UR-C465-haddi}



\begin{quote}
Before the new one: say the Urdu for an intestine, then the Urdu for a vein.
\end{quote}

\subsection*{You'll want to know: \ur{ہڈی}}
\textbf{\ur{ہڈی}} — \emph{haḍḍī} — \textquotedblleft{}a bone\textquotedblright{}.

\begin{quote}
\ur{کتا} \ur{ہڈی} \ur{کھا} \ur{رہا} \ur{تھا۔}
\end{quote}

\begin{quote}
\emph{kuttā haḍḍī khā rahā thā.}
\end{quote}

\begin{quote}
The dog was eating a bone.
\end{quote}

Feminine.

\begin{cousinweb}
\textbf{\ur{ہڈی}} (\emph{haḍḍī}) is inherited: CDIAL lists the Hindi cousin \emph{haḍḍī} under Sanskrit \emph{haḍḍa}, \textquotedblleft{}bone\textquotedblright{} (CDIAL 13952).
\end{cousinweb}

\begin{grammarlens}[title={What you've built}]
One more word for describing what you lived and felt.
\end{grammarlens}

\subsection*{Your turn}
\begin{itemize}
\raggedright
  \item \emph{Say it:} \emph{haḍḍī}
  \item \emph{Say it:} \emph{haḍḍī}, once more
  \item \emph{Say it:} \emph{kuttā haḍḍī khā rahā thā.} — The dog was eating a bone.
  \item \emph{Recall:} say the Urdu for an intestine, then the Urdu for a vein, then say \emph{haḍḍī} again
\end{itemize}

\begin{tcolorbox}[breakable,colback=teal!4,colframe=teal!35!black,title={Before you move on}]
What does \ur{ہڈی} mean? (\textquotedblleft{}A bone\textquotedblright{}.) Say it once more.
\end{tcolorbox}
\section[\texorpdfstring{\ur{پسلی} (paslī)}{پسلی (paslī)}]{\ur{پسلی} (paslī) — a rib}

\label{lesson:UR-C465-pasli}



\begin{quote}
Before the new one: say the Urdu for a vein, then the Urdu for a bone.
\end{quote}

\subsection*{You'll want to know: \ur{پسلی}}
\textbf{\ur{پسلی}} — \emph{paslī} — \textquotedblleft{}a rib\textquotedblright{}.

\begin{quote}
\ur{گرنے} \ur{سے} \ur{اس} \ur{کی} \ur{پسلی} \ur{میں} \ur{درد} \ur{ہے۔}
\end{quote}

\begin{quote}
\emph{girne se us kī paslī meṅ dard hai.}
\end{quote}

\begin{quote}
Her rib hurts since she fell.
\end{quote}

Feminine.

\begin{cousinweb}
\textbf{\ur{پسلی}} (\emph{paslī}) is inherited from Sanskrit \emph{párśu}, \textquotedblleft{}rib\textquotedblright{} (CDIAL 7948).
\end{cousinweb}

\begin{grammarlens}[title={What you've built}]
One more word for describing what you lived and felt.
\end{grammarlens}

\subsection*{Your turn}
\begin{itemize}
\raggedright
  \item \emph{Say it:} \emph{paslī}
  \item \emph{Say it:} \emph{paslī}, once more
  \item \emph{Say it:} \emph{girne se us kī paslī meṅ dard hai.} — Her rib hurts since she fell.
  \item \emph{Recall:} say the Urdu for a vein, then the Urdu for a bone, then say \emph{paslī} again
\end{itemize}

\begin{tcolorbox}[breakable,colback=teal!4,colframe=teal!35!black,title={Before you move on}]
What does \ur{پسلی} mean? (\textquotedblleft{}A rib\textquotedblright{}.) Say it once more.
\end{tcolorbox}
\section[\texorpdfstring{\ur{کھوپڑی} (khopṛī)}{کھوپڑی (khopṛī)}]{\ur{کھوپڑی} (khopṛī) — a skull}

\label{lesson:UR-C465-khopri}



\begin{quote}
Before the new one: say the Urdu for a bone, then the Urdu for a rib.
\end{quote}

\subsection*{You'll want to know: \ur{کھوپڑی}}
\textbf{\ur{کھوپڑی}} — \emph{khopṛī} — \textquotedblleft{}a skull\textquotedblright{}.

\begin{quote}
\ur{سر} \ur{کے} \ur{اندر} \ur{کھوپڑی} \ur{ہوتی} \ur{ہے۔}
\end{quote}

\begin{quote}
\emph{sar ke andar khopṛī hotī hai.}
\end{quote}

\begin{quote}
Inside the head is the skull.
\end{quote}

Feminine.

\begin{cousinweb}
\textbf{\ur{کھوپڑی}} (\emph{khopṛī}) is inherited: CDIAL lists the Hindi cousin \emph{khopṛī}, \textquotedblleft{}skull, shell\textquotedblright{}, under \emph{khōppa}, \textquotedblleft{}a coconut shell\textquotedblright{}, a form it reconstructs (CDIAL 3936).
\end{cousinweb}

\begin{grammarlens}[title={What you've built}]
One more word for describing what you lived and felt.
\end{grammarlens}

\subsection*{Your turn}
\begin{itemize}
\raggedright
  \item \emph{Say it:} \emph{khopṛī}
  \item \emph{Say it:} \emph{khopṛī}, once more
  \item \emph{Say it:} \emph{sar ke andar khopṛī hotī hai.} — Inside the head is the skull.
  \item \emph{Recall:} say the Urdu for a bone, then the Urdu for a rib, then say \emph{khopṛī} again
\end{itemize}

\begin{tcolorbox}[breakable,colback=teal!4,colframe=teal!35!black,title={Before you move on}]
What does \ur{کھوپڑی} mean? (\textquotedblleft{}A skull\textquotedblright{}.) Say it once more.
\end{tcolorbox}
\section[\texorpdfstring{\ur{ریڑھ} (rīṛh)}{ریڑھ (rīṛh)}]{\ur{ریڑھ} (rīṛh) — the backbone}

\label{lesson:UR-C465-rirh}



\begin{quote}
Before the new one: say the Urdu for a rib, then the Urdu for a skull.
\end{quote}

\subsection*{You'll want to know: \ur{ریڑھ}}
\textbf{\ur{ریڑھ}} — \emph{rīṛh} — \textquotedblleft{}the backbone\textquotedblright{}.

\begin{quote}
\ur{بیٹھتے} \ur{وقت} \ur{ریڑھ} \ur{کی} \ur{ہڈی} \ur{سیدھی} \ur{رکھیے۔}
\end{quote}

\begin{quote}
\emph{baiṭhte vaqt rīṛh kī haḍḍī sīdhī rakhiye.}
\end{quote}

\begin{quote}
Keep your back straight when you sit.
\end{quote}

Feminine.

\begin{cousinweb}
\textbf{\ur{ریڑھ}} (\emph{rīṛh}) is inherited from Sanskrit \emph{rīḍhaka}, \textquotedblleft{}backbone\textquotedblright{} (CDIAL 10749a).
\end{cousinweb}

\begin{grammarlens}[title={What you've built}]
One more word for describing what you lived and felt.
\end{grammarlens}

\subsection*{Your turn}
\begin{itemize}
\raggedright
  \item \emph{Say it:} \emph{rīṛh}
  \item \emph{Say it:} \emph{rīṛh}, once more
  \item \emph{Say it:} \emph{baiṭhte vaqt rīṛh kī haḍḍī sīdhī rakhiye.} — Keep your back straight when you sit.
  \item \emph{Recall:} say the Urdu for a rib, then the Urdu for a skull, then say \emph{rīṛh} again
\end{itemize}

\begin{tcolorbox}[breakable,colback=teal!4,colframe=teal!35!black,title={Before you move on}]
What does \ur{ریڑھ} mean? (\textquotedblleft{}The backbone\textquotedblright{}.) Say it once more.
\end{tcolorbox}
\section[\texorpdfstring{\ur{جبڑا} (jabṛā)}{جبڑا (jabṛā)}]{\ur{جبڑا} (jabṛā) — a jaw}

\label{lesson:UR-C465-jabra}



\begin{quote}
Before the new one: say the Urdu for a skull, then the Urdu for the backbone.
\end{quote}

\subsection*{You'll want to know: \ur{جبڑا}}
\textbf{\ur{جبڑا}} — \emph{jabṛā} — \textquotedblleft{}a jaw\textquotedblright{}.

\begin{quote}
\ur{شیر} \ur{کا} \ur{جبڑا} \ur{بڑا} \ur{مضبوط} \ur{ہوتا} \ur{ہے۔}
\end{quote}

\begin{quote}
\emph{sher kā jabṛā baṛā mazbūt hotā hai.}
\end{quote}

\begin{quote}
A lion's jaw is very strong.
\end{quote}

Masculine.

\begin{cousinweb}
\textbf{\ur{جبڑا}} (\emph{jabṛā}) is inherited: CDIAL lists the Hindi cousin \emph{jābṛā}, \textquotedblleft{}jaw-bone\textquotedblright{}, under Sanskrit \emph{jámbha}, \textquotedblleft{}tooth, jaw\textquotedblright{} (CDIAL 5137).
\end{cousinweb}

\begin{grammarlens}[title={What you've built}]
One more word for describing what you lived and felt.
\end{grammarlens}

\subsection*{Your turn}
\begin{itemize}
\raggedright
  \item \emph{Say it:} \emph{jabṛā}
  \item \emph{Say it:} \emph{jabṛā}, once more
  \item \emph{Say it:} \emph{sher kā jabṛā baṛā mazbūt hotā hai.} — A lion's jaw is very strong.
  \item \emph{Recall:} say the Urdu for a skull, then the Urdu for the backbone, then say \emph{jabṛā} again
\end{itemize}

\begin{tcolorbox}[breakable,colback=teal!4,colframe=teal!35!black,title={Before you move on}]
What does \ur{جبڑا} mean? (\textquotedblleft{}A jaw\textquotedblright{}.) Say it once more.
\end{tcolorbox}
